% ============================================================
%  Unit 6 (continued) — Collisions & COM Frame
% ============================================================

\chapter{Collisions and Conservation Laws}

\section{Two-Body Collisions}

In a two-body collision momentum is always conserved.
\begin{itemize}
    \item \textbf{Elastic collision:} energy is conserved.
    \item If blocks stick together, the collision is \textbf{inelastic}.
\end{itemize}

\subsection{Velocities in the COM Frame}

Let $v_{1,\text{com}}$ and $v_{2,\text{com}}$ be velocities in the COM frame. Then
\begin{align*}
    \vect{v}_{1,\text{com}} &= \vect{v}_1 - \vect{v}_\text{com}
        = \frac{m_2}{m_1+m_2}(\vect{v}_1 - \vect{v}_2) \\
    \vect{v}_{2,\text{com}} &= \vect{v}_2 - \vect{v}_\text{com}
        = \frac{m_1}{m_1+m_2}(\vect{v}_1 - \vect{v}_2)
\end{align*}

The momenta in the COM frame are
\begin{align*}
    \vect{p}_{1,\text{com}} &= m_1 \vect{v}_{1,\text{com}}
        = \frac{m_1 m_2}{m_1+m_2}(\vect{v}_1 - \vect{v}_2) = \mu\,\Delta v \\
    \vect{p}_{2,\text{com}} &= m_2 \vect{v}_{2,\text{com}}
        = \frac{m_2 m_1}{m_1+m_2}(\vect{v}_2 - \vect{v}_1)
\end{align*}

Thus,
\[
    \boxed{\vect{p}_{1,\text{com}} = -\vect{p}_{2,\text{com}}}
\]

because the COM is at the origin — \textbf{total momentum is 0}.

\[
    \vect{p}_{\text{com,net}} = \vec{0}
\]

The total kinetic energy in the COM frame is $\tfrac{1}{2}m v_{1,\text{com}}^2 + \tfrac{1}{2}m v_{2,\text{com}}^2$.

\begin{fact}[Elastic collisions in the COM frame]
For elastic collisions,
\[
    \bigl|\vect{v}_{1,\text{com}}\bigr|_{\text{before}}
    = \bigl|\vect{v}_{1,\text{com}}\bigr|_{\text{after}}
\]
Speeds are unchanged by the collision (useful fact — prove this fact).
\end{fact}

\subsection{Elastic Collision Example — Collision of the Carts}

\begin{example}[Elastic collision of two carts]
Initial velocities $\vect{v}_1,\,\vect{v}_2$; final velocities $\vect{v}_1',\,\vect{v}_2'$.

Derive $v_1 + v_1' = v_2 + v_2'$.

Conservation of momentum (projection in 1D):
\[
    m_1 v_1 + m_2 v_2 = m_1 v_1' + m_2 v_2'
\]
Conservation of energy:
\[
    m_1 v_1^2 + m_2 v_2^2 = m_1 {v_1'}^2 + m_2 {v_2'}^2
\]
\begin{align*}
    m_1(v_1^2-{v_1'}^2) &= m_2({v_2'}^2-v_2^2) \\
    m_1(v_1-v_1')(v_1+v_1') &= m_2(v_2'-v_2)(v_2'+v_2) \\
    m_1(v_1-v_1') &= m_2(v_2'-v_2)
\end{align*}

Thus the \textbf{useful collision fact}:
\[
    \boxed{v_1 + v_1' = v_2 + v_2'}
\]
\end{example}

\subsection{Inelastic Collision}

Energy is \textbf{NOT} conserved.

\textbf{Completely inelastic collision} $\Rightarrow$ carts stick together.

\[
    m_1\vect{v}_1 + m_2\vect{v}_2 = (m_1+m_2)\vect{v}_f'
\]

\subsubsection*{Projection in 1D}
\[
    v_f = \frac{m_1 v_1 + m_2 v_2}{m_1 + m_2}
\]

\subsection{Elastic Scattering in the COM Frame — Example}

\begin{example}[Elastic scattering through angle $\theta$ in COM frame]
A particle of mass $m$ and speed $v_0$ collides elastically with a particle of mass $M$ which is initially at rest. Particle $m$ is scattered through an angle $\theta$ in the COM frame. Find the final speed in the lab frame.

\subsubsection*{Lab frame (before)}
$m$ moving at $v_0$, $M$ at rest.

\subsubsection*{COM frame (before)}
\begin{align*}
    \vect{v}_{m,\text{com}} &= \vect{v}_0 - \vect{v}_\text{com} \\
    |\vect{v}_{m,\text{com}}| &= \frac{Mv_0}{m+M} \\
    \vect{v}_{M,\text{com}} &= -\vect{v}_\text{com} = -\frac{mv_0}{m+M}
\end{align*}

By the useful fact (speeds unchanged in COM frame under elastic collision):
\[
    v_{m,\text{com}}' = \frac{Mv_0}{m+M}
\]

\subsubsection*{Going back to the lab frame}
\[
    \vect{v}_m' = \vect{v}_{1,\text{com}}' + \vect{v}_\text{com}
\]

\begin{figure}[H]
\centering
\caption{Vector diagram for returning to lab frame: the final lab velocity $\vect{v}_m'$ is the vector sum of the COM-frame velocity $\vect{v}_{1,\text{com}}'$ and the COM velocity $\vect{v}_\text{com}$.}
\end{figure}
\end{example}

\subsection{Bullet Hitting a Pendulum}

\begin{figure}[H]
\begin{center}
\begin{tikzpicture}[scale=0.9, >=Stealth]
  % Pivot
  \fill[pattern=north east lines] (-0.4,3.0) rectangle (0.4,3.3);
  \draw[thick] (-0.4,3.0) -- (0.4,3.0);
  \fill (0,3.0) circle[radius=0.05];
  % String (at rest, vertical)
  \draw[thick] (0,3.0) -- (0,1.2);
  % Bob
  \fill[gray!50] (0,1.0) circle[radius=0.2];
  \draw[thick] (0,1.0) circle[radius=0.2];
  \node[right] at (0.22,1.0) {$M$};
  % String length
  \node[left] at (-0.1,2.1) {$\ell$};
  % Bullet
  \fill[red!70] (-1.5,1.0) -- (-0.9,0.88) -- (-0.9,1.12) -- cycle;
  \draw[->,thick,red] (-1.8,1.0) -- (-0.25,1.0) node[above,pos=0.3]{$m,\;v_i$};
  % After collision: swung to angle phi
  \draw[dashed,gray] (0,3.0) -- (0,0.8);
  \draw[thick,blue] (0,3.0) -- (0.85,1.55);
  \fill[blue!50] (0.85,1.35) circle[radius=0.2];
  \draw[thick,blue] (0.85,1.35) circle[radius=0.2];
  % Angle arc
  \draw (0,2.0) arc[start angle=270,end angle=242,radius=1.0];
  \node[below left] at (0.1,2.0) {$\phi$};
  % Height h
  \draw[<->,gray] (1.3,1.35) -- (1.3,1.0) node[midway,right]{$h$};
\end{tikzpicture}
\end{center}
\caption{A bullet of mass $m$ moving at $v_i$ embeds in a pendulum bob of mass $M$. After the inelastic collision, the combined mass $m+M$ swings to angle $\phi$ and height $h = \ell(1-\cos\phi)$.}
\label{fig:bullet_pendulum}
\end{figure}

From momentum conservation (during collision) and energy conservation (during swing):
\[
    \boxed{v_i = \frac{(m+M)}{m}\sqrt{2g\ell(1-\cos\phi)}}
\]

% ============================================================
%  Chapter 7 — Angular Momentum and Rotation
% ============================================================

\chapter{Angular Momentum and Rotation}

\section{Angular Momentum}

\begin{definition}[Angular Momentum]
\[
    \vect{L} = \vect{r} \times \vect{p}
\]
where $\vect{r}$ is the position of the object at a point and $\vect{p}$ is its momentum.
\end{definition}

Angular momentum:
\begin{itemize}
    \item Describes the motion of an object around a point.
    \item Tells us the rotation of the system.
\end{itemize}

For a particle in circular motion (rotation confined to a plane):
\begin{align*}
    \vect{L} &= \vect{r} \times \vect{p} = \vect{r} \times m\vect{v} = m(\vect{r}\times\vect{v})
\end{align*}

\subsubsection*{Computing $\vect{L}_A$ (pivot at centre of orbit)}

\begin{align*}
    \vect{L}_A &= m(\vect{r}_A \times \vect{v}) \\
    &= m\bigl(r_A\,\hat{r} \times v\,\hat{\theta}\bigr) \\
    &= mr_A v\,(\hat{r}\times\hat{\theta}) \\
    &= mr_A v\,\hat{z}
\end{align*}

\subsubsection*{Computing $\vect{L}_B$ (pivot off-axis)}

\begin{align*}
    \vect{L}_B &= m\bigl((-x\hat{z}+r_A\hat{r})\times v\hat{\theta}\bigr) \\
    &= m\bigl(-xv(\hat{z}\times\hat{\theta}) + r_A v(\hat{r}\times\hat{\theta})\bigr) \\
    &= m\bigl(-xv(-\hat{r}) + r_A v\,\hat{z}\bigr) \\
    &= mxv\,\hat{r} + mr_A v\,\hat{z}
\end{align*}

\subsection{DART Mission (Example)}

\begin{figure}[H]
\centering
\caption{DART spacecraft ($11\times10^6\,\text{km}$ away, mass $M=610\,\text{kg}$, 90\,s transmission delay) impacting Dimorphos (diameter $\approx163\,\text{m}$, rock density $2\,\text{g/cm}^3$).}
\end{figure}

\textbf{Cross-section:} area that collision sees:
\[
    A_c = \pi r^2
\]

\[
    \tan\theta \approx \frac{80\,\text{m}}{11\times10^9} \approx 7\times10^{-9}\,\text{rad}
    \implies \text{not possible to aim with that precision}
\]

\textbf{However, gravity works in our favour.}

$b$ — impact parameter (largest impact parameter that would still hit the asteroid):
\[
    A_g \pi R^2 = \pi b^2 \qquad \text{(if no gravity, }B'=R,\;A_g=1\text{)}
\]
\[
    A_g > 1
\]

Goal: find $b'$ (before).
\begin{enumerate}
    \item Energy conservation — True
    \item Angular momentum conservation — True
    \item Momentum conservation — No (external force: gravity)
\end{enumerate}

\section{Chasles' Theorem}

\begin{theorem}[Chasles' Theorem]
Any displacement is the combination of a single rotation and a translation.
\end{theorem}

\section{Angular Momentum of a Particle}

\[
    \vect{L} = \vect{r}\times\vect{p}
\]

where $\vect{r}$ is position with respect to a given coordinate system and $\vect{p}$ is the momentum vector of the particle.

\[
    |\vect{L}| = |\vect{r}||\vect{p}|\sin\alpha
\]

We can break $\vect{r}$ into its components: $r_\perp$ (perpendicular to trajectory) and $r_\parallel$ (parallel to trajectory).

Thus,
\[
    r_\perp = r\sin(\pi-\phi) = r\sin\phi
\]
\[
    L = r_\perp\, p \sin\phi = r_\perp\, p
\]

Or: $L = r\, p_\perp$.

\textbf{In summary,}
\[
    |L_z| = |\vect{r}||\vect{p}| = |\vect{r}||\vect{p}_\perp|
\]

(where $\vect{r}\cdot\vect{p}=0$) and $\vec{L}_z = m|\vect{r}|^2\omega\,\hat{k}$.

If motion is 2D:
\[
    \boxed{\vect{L} = \vect{r}\times\vect{p} = m(xv_y - yv_x)\,\hat{k}}
\]

\begin{example}[Conical pendulum angular momentum]
Because $\vect{r}\perp\vect{p}$,
\[
    |\vect{L}_A| = |\vect{r}||\vect{p}| = rp
\]

Recall $|p| = Mv = Mr\omega$, thus
\[
    \boxed{\vect{L}_A = Mr^2\omega\,\hat{k}}
    \quad\text{(when position and momentum are perpendicular)}
\]

Now imagine momentum with respect to pivot at point $B$:
\[
    |\vect{L}_B| = |\vect{r}'||\vect{p}| = m|\ell||\vect{v}| = \ell |p|
    = M\ell r\omega
\]

The direction of $\vect{L}_B$ is \emph{not} constant (connects to the binormal vector).
\end{example}

\section{Fixed-Axis Rotation}

\begin{definition}[Fixed-axis rotation]
The axis can \underline{translate} but \underline{cannot} rotate.
\end{definition}

\begin{itemize}
    \item $\rho$ — perpendicular distance to the axis of rotation.
    \item $r$ — distance from the origin $= \sqrt{x^2+y^2+z^2}$.
\end{itemize}

For the $j$-th particle undergoing circular motion about the $z$-axis,
\[
    |\vect{v}_j| = |\dot{\vect{r}}_j| = \omega\rho_j \tag{1}
\]

where $\rho_j = \sqrt{x^2+y^2}$ in this case.

Thus, the angular momentum of the $j$-th particle is
\[
    \vect{L}_j = \vect{r}_j \times m_j \vect{v}_j
\]

We only care about angular momentum along the axis of rotation, so
\[
    L_{j,z} = m_j v_j \times \rho_j = m_j v_j \rho_j
\]
\[
    \Rightarrow \quad L_{j,z} = m_j \rho_j^2 \omega
\]

Thus, total angular momentum:
\[
    L_z = \sum_j L_{j,z} = \sum_j m_j \rho_j^2\,\omega \tag{2}
\]

\section{Moment of Inertia}

Rewriting equation (2),
\[
    L_z = \omega\!\left(\sum_j m_j \rho_j^2\right) = I\omega
\]

\begin{definition}[Moment of Inertia]
\[
    \boxed{I = \sum_j m_j \rho_j^2}
\]
\end{definition}

For a continuous distribution of matter,
\[
    I = \int \rho^2\,dm = \int(x^2+y^2)\,dm = \int(x^2+y^2)\,w\,dV
\]
where $w$ is the mass per unit volume (density).

\subsection{Moment of Inertia of Common Objects}

\subsubsection{(1) Uniform Thin Ring}

\begin{figure}[H]
\begin{center}
\begin{tikzpicture}[scale=0.9, >=Stealth]
  % Ring
  \draw[thick] (0,0) circle[radius=1.2];
  \draw[dashed,gray] (0,0) circle[radius=0.08];
  \fill (0,0) circle[radius=0.05];
  % Radius label
  \draw[->,gray] (0,0) -- (0.85,0.85) node[above right,scale=0.85]{$R$};
  % Element dm
  \draw[thick,red,fill=red!30] (1.1,0.42) arc[start angle=21,end angle=30,radius=1.2]
    -- (0.93,0.58) arc[start angle=32,end angle=22,radius=1.0] -- cycle;
  \node[red,scale=0.85] at (1.45,0.6) {$dm$};
  % Rotation axis (z)
  \draw[->] (0,-1.6) -- (0,-0.15) node[right,scale=0.8]{$z$-axis};
  % Label
  \node[below] at (0,-1.8) {$I = MR^2$};
\end{tikzpicture}
\end{center}
\caption{Uniform thin ring of mass $M$ and radius $R$ rotating about its central axis. Every mass element is at the same distance $R$ from the axis, giving $I = MR^2$.}
\label{fig:moi_ring}
\end{figure}

\[
    I = \int \rho^2\,dm = \int_0^{2\pi R} R^2 \cdot \frac{M}{2\pi R}\,ds
    = \frac{M}{2\pi} R \Bigl[s\Bigr]_0^{2\pi R} = MR^2
\]

\[
    \boxed{I = MR^2}
\]

\subsubsection{(2) Uniform Thin Disk}

Using double integration with surface mass density $\sigma = M/(\pi R^2)$ (mass per unit area):

\[
    I = \int \rho^2\,dm = \int \rho^2\,\sigma\,dA
\]

In polar coordinates, $dA = \rho\,d\rho\,d\theta$, so:

\[
    I = \frac{M}{\pi R^2}\int_0^{2\pi}\int_0^{R} \rho^3\,d\rho\,d\theta
    = \frac{M}{\pi R^2}\int_0^{2\pi}\frac{R^4}{4}\,d\theta
    = \frac{M}{\pi R^2}\cdot\frac{2\pi R^4}{4} = \frac{1}{2}MR^2
\]

\[
    \boxed{I = \tfrac{1}{2}MR^2}
\]

\subsubsection{(3) Uniform Thin Stick (about centre)}

\begin{figure}[H]
\begin{center}
\begin{tikzpicture}[scale=0.9, >=Stealth]
  % Stick
  \draw[thick,gray!70,fill=gray!30] (-2.0,-0.12) rectangle (2.0,0.12);
  % Centre marker
  \fill (0,0) circle[radius=0.06];
  \draw[->,gray,dashed] (0,-0.5) -- (0,0.5) node[above,scale=0.8]{axis};
  % Length arrows
  \draw[<->] (-2.0,-0.4) -- (2.0,-0.4) node[midway,below]{$L$};
  % Element dx
  \draw[thick,red,fill=red!30] (0.6,-0.12) rectangle (0.85,0.12);
  \node[red,above,scale=0.8] at (0.72,0.12) {$dx$};
  \draw[<->] (0,-0.65) -- (0.72,-0.65) node[midway,below,scale=0.8]{$x$};
  % Label
  \node[below] at (0,-0.9) {$I = \tfrac{1}{12}ML^2$};
\end{tikzpicture}
\end{center}
\caption{Uniform thin stick of mass $M$ and length $L$ rotating about its centre ($x=0$). A mass element $dm = (M/L)\,dx$ at position $x$ contributes $x^2\,dm$ to the moment of inertia.}
\label{fig:moi_stick}
\end{figure}

With $dm = \frac{M}{L}\,dx = \lambda\,dx$:

\[
    I = \int x^2\,dm = \lambda\int_{-L/2}^{L/2} x^2\,dx
    = \frac{M}{L}\cdot\frac{x^3}{3}\Bigg|_{-L/2}^{L/2}
    = \frac{M}{L}\!\left(\frac{L^3}{24}+\frac{L^3}{24}\right)
    = \frac{1}{12}ML^2
\]

\[
    \boxed{I = \tfrac{1}{12}ML^2}
\]

\subsubsection{(4) Uniform Thin Stick (about perpendicular end axis)}

\[
    I_\text{rod} = \frac{M}{L}\int_0^{L} x^2\,dx = \boxed{\tfrac{1}{3}ML^2}
\]

\subsubsection{(5) Uniform Sphere}

\[
    I_\text{sphere} = \tfrac{2}{5}MR^2
\]

\section{Parallel Axis Theorem}

\begin{theorem}[Parallel Axis Theorem]
\[
    I = I_0 + M\ell^2
\]
where $I_0$ is the moment of inertia about an axis through the COM, and $\ell$ is the distance between the old and new axis.
\end{theorem}

\begin{figure}[H]
\begin{center}
\begin{tikzpicture}[scale=0.85, >=Stealth]
  % Object (ellipse)
  \draw[thick,gray!60,fill=gray!20] (0,0) ellipse[x radius=1.5, y radius=0.8];
  % COM axis
  \fill (0,0) circle[radius=0.06];
  \draw[dashed,blue,thick] (0,-1.2) -- (0,1.2) node[above,blue,scale=0.8]{COM axis};
  % New axis
  \fill (1.5,0) circle[radius=0.06];
  \draw[dashed,red,thick] (1.5,-1.2) -- (1.5,1.2) node[above,red,scale=0.8]{new axis};
  % Distance ell
  \draw[<->,ForestGreen,thick] (0,-1.0) -- (1.5,-1.0) node[midway,below]{$\ell$};
  % Label
  \node[below] at (0.75,-1.4) {$I = I_0 + M\ell^2$};
\end{tikzpicture}
\end{center}
\caption{Parallel axis theorem: shifting the rotation axis from the COM (blue) to a parallel axis at distance $\ell$ (red) increases the moment of inertia by $M\ell^2$.}
\label{fig:parallel_axis}
\end{figure}

\subsubsection*{Proof}

Choose $I_0$ to lie along the $z$-axis. Then
\[
    \vec{\rho}_j = x_j\,\ihat + y_j\,\jhat, \qquad I = \sum_j m_j \rho_j^2
\]

From the diagram, $\vec{\rho}_j = \vec{R}_\perp + \vec{\rho}_j'$, so:
\begin{align*}
    I &= \sum_j m_j \rho_j^2 = \sum_j m_j\bigl(\vec{R}_\perp+\vec{\rho}_j'\bigr)^2 \\
    &= \sum_j m_j\bigl(R_\perp^2 + 2\vec{\rho}_j'\cdot\vec{R}_\perp + {\rho_j'}^2\bigr) \\
    &= \sum_j m_j R_\perp^2 + \sum_j m_j(2\vec{\rho}_j'\cdot\vec{R}_\perp)
       + \sum_j m_j {\rho_j'}^2
\end{align*}

The middle term $= \sum_j(2\vec{\rho}_j' m_j \cdot \vec{R}_\perp) = 0$, because $\vec{\rho}_j'$ is measured with respect to the COM.

Therefore,
\[
    I = MR_\perp^2 + I_0 \qquad (\text{where }R_\perp \equiv \ell)
\]

\section{Torque}

\begin{definition}[Torque]
\[
    \vec{\tau} = \vect{r}\times\vect{F}
\]
Torque is the angular analogue to force.
\end{definition}

As usual,
\[
    |\vec{\tau}| = |r_\perp||\vect{F}| = |\vect{r}||F_\perp|
\]

\subsection{Example: Torque Due to Gravity}

Find the torque in a uniform gravitational field where $\vect{F}_g = m\vect{g}$.

\begin{align*}
    \vec{\tau}_j &= \vect{r}_j \times m_j \vect{g} \\
    \vec{\tau}_\text{net} &= \sum_j \vec{\tau}_j
        = \sum_j \vect{r}_j \times m_j \vect{g}
        = \sum_j m_j \vect{r}_j \times \vect{g}
        = \left(\sum_j m_j \vect{r}_j\right)\times\vect{g} \\
    &= M\vec{R}_\text{com}\times\vect{g} = \vec{R}\times\vec{W}
\end{align*}

\textbf{Torque is 0 at equilibrium} (and $\vect{F}_\text{net}=\vec{0}$).

\section{Torque and Angular Momentum}

\begin{formula}[Newton's 2nd law for rotation]
\[
    \boxed{\dv{\vect{L}}{t}} = \dv{}{t}(\vect{r}\times\vect{p})
    = \dot{\vect{r}}\times\vect{p} + \vect{r}\times\dot{\vect{p}}
\]

The first term $\dot{\vect{r}}\times\vect{p} = \vect{v}\times m\vect{v} = 0$.

Thus,
\[
    \dv{\vect{L}}{t} = \vect{r}\times\dot{\vect{p}} = \vect{r}\times\vect{F} = \vec{\tau}
\]
\end{formula}

\section{Conservation of Angular Momentum}

\begin{theorem}[Conservation of Angular Momentum]
If $\vec{\tau} = \vec{0}$, then $\vect{L}$ is constant.
\end{theorem}

\subsection{Example: Central Force Motion and Equal Areas}

\begin{example}[Kepler's Second Law from angular momentum conservation]
Consider a particle moving under a central force
\[
    \vect{F}(r) = f(r)\,\hat{r}
\]
Then
\[
    \vec{\tau} = \vect{r}\times f(r)\hat{r} = 0
    \quad\text{(both pointing in the same direction)}
\]
Thus $\vect{L}$ is conserved (equal areas swept in equal times).

\subsubsection*{Rate of area swept}

Area swept out:
\[
    \Delta A \approx \tfrac{1}{2}(r+\Delta r)\,r\,\Delta\theta
    = \tfrac{1}{2}r^2\,\Delta\theta + \tfrac{1}{2}\Delta r\cdot r\,\Delta\theta
\]
\[
    \dv{A}{t} = \lim\left(\frac{\Delta A}{\Delta t}\right)
    = \frac{1}{2}\!\left(r^2\frac{\Delta\theta}{\Delta t} + r\frac{\Delta\theta\,\Delta r}{\Delta t}\right)
    = \frac{1}{2}r^2\dv{\theta}{t}
\]

From polar coordinates, $\vect{v} = \dot{r}\hat{r} + r\dot{\theta}\hat{\theta}$, so
\begin{align*}
    \vect{L} &= \vect{r}\times m\vect{v}
        = r\hat{r}\times m(\dot{r}\hat{r}+r\dot{\theta}\hat{\theta}) \\
    &= r\hat{r}\times m\dot{r}\hat{r} + r\hat{r}\times mr\dot{\theta}\hat{\theta} \\
    &= mr^2\dot{\theta}\,\hat{z} = 2m\dv{A}{t} = C
\end{align*}

Thus,
\[
    \boxed{\dv{A}{t} = C}
\]
\end{example}

\subsection{Example: Capture Cross-Section of a Planet}

\begin{example}[Gravitational capture cross-section]
Find the effective area $A_e$ that it takes to hit a planet.

\begin{figure}[H]
\centering
\caption{A projectile with impact parameter $b$ (perpendicular distance from the asymptotic trajectory to the centre), moving with initial speed $v_0$ from far away, is gravitationally focused toward a planet of radius $R$. The effective cross-section $A_e = \pi(b')^2$ is larger than the geometric cross-section $A_g = \pi R^2$.}
\end{figure}

$b$ is the \underline{impact parameter}; $\vec{a}$ is parallel to the initial trajectory.

Assumed that the distance from launch point to planet is extremely large compared to $R$, so that $b_1,b_1',b_2$ are effectively parallel.

Thus, $A_e = \pi(b')^2$ (not linear).

Energy and angular momentum are conserved.

As $r\to\infty$:
\[
    E_i = \tfrac{1}{2}mv_0^2, \qquad L_i = -mb'v_0
\]
(rotation is clockwise ($-$) and $b'$ is the perpendicular component of $\vect{r}$).

Collision occurs when the point of closest approach is $R$; at this point $\vect{r}\perp\vect{v}$.

If speed is $v(R)$ at closest approach:
\[
    L_C = -mR\,v(R), \qquad E_C = \tfrac{1}{2}mv(R)^2 - \frac{GmM}{R}
\]

From conservation of angular momentum:
\[
    -mb'v_0 = -mR\,v(R) \tag{1}
\]

From conservation of energy:
\[
    \tfrac{1}{2}mv_0^2 = \tfrac{1}{2}mv(R)^2 - \frac{GmM}{R} \tag{2}
\]

(1) gives $v(R) = \dfrac{b'v_0}{R}$.

(2) gives
\[
    (b')^2 = R^2\!\left(1+\frac{GmM/R}{\frac{1}{2}mv_0^2}\right)
    = R^2\!\left(1+\frac{2GM/R}{v_0^2}\right)
\]

Therefore,
\[
    A_e = \pi(b')^2 = \pi R^2\!\left(1+\frac{2GM/R}{\frac{1}{2}mv_0^2/\frac{1}{2}m}\right)
    > \pi R^2
\]
\[
    = A_g\!\left(1 - \frac{U(R)}{E}\right)
\]

Checks:
\begin{itemize}
    \item If $U(R)\to 0$, then $A_e \to A_g$. \checkmark
    \item As $E\to 0$, $A_e\to\infty$. \checkmark
\end{itemize}
\end{example}

\section{Dynamics of a Conical Pendulum}

\begin{example}[Conical pendulum]
Bob undergoes uniform circular motion. Let $\alpha$ be the half-angle, $\ell = |r_B|$ the string length, and $\omega$ the angular velocity.

Equations of motion:
\[
    T\cos\alpha - mg = 0
\]
\[
    \vect{F} = -T\sin\alpha\,\hat{r}
\]
\end{example}
